% Generated draft: requires mathematical and source review.
\documentclass[11pt]{article}
\usepackage[margin=25mm]{geometry}
\usepackage{fontspec}
\setmainfont{Latin Modern Roman}
\usepackage{amsmath,amssymb,mathtools,unicode-math}
\setmathfont{Latin Modern Math}
\usepackage{enumitem,needspace,xurl,hyperref}
\setlength{\parindent}{0pt}
\setlength{\parskip}{5pt}
\setlength{\emergencystretch}{3em}
\newcommand{\N}{\mathbb{N}}
\newcommand{\Z}{\mathbb{Z}}
\newcommand{\Q}{\mathbb{Q}}
\newcommand{\R}{\mathbb{R}}
\newcommand{\C}{\mathbb{C}}
\newcommand{\F}{\mathbb{F}}
\newcommand{\A}{\mathbb{A}}
\newcommand{\PP}{\mathbb{P}}
\newcommand{\E}{\mathbb{E}}
\newcommand{\eps}{\varepsilon}
\newcommand{\dd}{\,\mathrm{d}}
\begin{document}
\section*{UnsolvedMath ID: 800126}
\textbf{Source problem code:} catalogue entry
\subsection{Definitions and mathematical statement}
A PL manifold is a manifold equipped with a piecewise linear structure. A PL homeomorphism is a homeomorphism that is piecewise linear with respect to the given PL structures. For a PL manifold \(M\), the source writes
\[
\mathrm{PL}_0(M)=\{h\colon M\to M:\ h\text{ is a PL homeomorphism isotopic to }\operatorname{id}_M\},
\]
with composition as the group operation. An isotopy is a continuous family of homeomorphisms joining the two specified maps. The source does not specify whether the intermediate homeomorphisms must be PL or whether the isotopy must itself satisfy a PL condition; these conventions remain unresolved here.

\textbf{Question (Epstein; as recorded by Nariman and Yazdi).} For every PL manifold \(M\) in the intended class, is \(\mathrm{PL}_0(M)\) a simple group? Explicitly, after fixing the intended manifold and isotopy conventions, does
\[
N\trianglelefteq\mathrm{PL}_0(M)
\quad\Longrightarrow\quad
N=\{\operatorname{id}_M\}\ \text{or}\ N=\mathrm{PL}_0(M)
\]
hold for every normal subgroup \(N\)? Under the convention that a simple group must be nontrivial, one must additionally require \(\mathrm{PL}_0(M)\ne\{\operatorname{id}_M\}\).

The supplied formulation does not restrict dimension, connectedness, compactness, or the presence of boundary. It also does not specify any compact support condition, boundary fixing condition, or corresponding condition on isotopies. These omissions are unresolved qualifications of the question; no such hypotheses are imposed in this draft.
\subsection{Short English statement}
Epstein's question, recorded by Nariman and Yazdi, asks whether the group of PL homeomorphisms of a PL manifold that are isotopic to the identity is simple. The intended class of manifolds and the isotopy, support, and boundary conventions require clarification. The supplied registry labels the question open; that label is not independently verified here.
\subsection{Sources}
\begin{itemize}
\item Local UnsolvedMath registry, supplied complete source record: permanent ID \(800126\), exact source problem code \texttt{catalogue entry}, title \emph{Foliations and Diffeomorphism Groups --- Question 67}. The record attributes the compilation to Nariman and Yazdi, gives proposed year 2025, and labels the problem open. Its creation and update timestamps are both \texttt{2026-10-16T14:00:00Z}; these are supplied registry metadata, not independent mathematical verification.
\item Nariman and Yazdi, \emph{Problem List on Foliations and Diffeomorphism Groups} (2025), Question 67, \url{https://arxiv.org/abs/2503.18274}. The supplied provenance records access on 2026-10-16 and extraction by \texttt{arxiv\_latex}; no source visit or verification is asserted here.
\item The supplied statement attributes the question to Epstein and cites the reference key \texttt{epstein1970simplicity}. No bibliographic entry or URL for that reference is supplied, so its exact identity and formulation remain unverified.
\end{itemize}
\end{document}
